\documentclass[11pt]{article}
\usepackage[utf8]{inputenc}
\usepackage{amsmath, amssymb}
\usepackage{xcolor}
\usepackage[margin=1in]{geometry}

\title{Technical Communication Assignment-8}
\author{UTKARSH KUMAR 23114101}
\date{April 2026}

\begin{document}
\maketitle

\section*{Question 1}

\subsection*{(a) Recursion Tree}
The recursion tree for \texttt{fun(arr, 0, 7)} is:

\begin{itemize}
    \item \texttt{fun(arr, 0, 7)} \quad [\texttt{mid = 3}]
    \begin{itemize}
        \item \texttt{fun(arr, 0, 3)} \quad [\texttt{mid = 1}]
        \begin{itemize}
            \item \texttt{fun(arr, 0, 1)} \quad [\texttt{mid = 0}]
            \begin{itemize}
                \item \texttt{fun(arr, 0, 0)} \quad $\to$ Base case
                \item \texttt{fun(arr, 1, 1)} \quad $\to$ Base case
            \end{itemize}
            \item \texttt{fun(arr, 2, 3)} \quad [\texttt{mid = 2}]
            \begin{itemize}
                \item \texttt{fun(arr, 2, 2)} \quad $\to$ Base case
                \item \texttt{fun(arr, 3, 3)} \quad $\to$ Base case
            \end{itemize}
        \end{itemize}
        \item \texttt{fun(arr, 4, 7)} \quad [\texttt{mid = 5}]
        \begin{itemize}
            \item \texttt{fun(arr, 4, 5)} \quad [\texttt{mid = 4}]
            \begin{itemize}
                \item \texttt{fun(arr, 4, 4)} \quad $\to$ Base case
                \item \texttt{fun(arr, 5, 5)} \quad $\to$ Base case
            \end{itemize}
            \item \texttt{fun(arr, 6, 7)} \quad [\texttt{mid = 6}]
            \begin{itemize}
                \item \texttt{fun(arr, 6, 6)} \quad $\to$ Base case
                \item \texttt{fun(arr, 7, 7)} \quad $\to$ Base case
            \end{itemize}
        \end{itemize}
    \end{itemize}
\end{itemize}

\subsection*{(b) Intermediate Arrays}
The intermediate arrays after each call:
\[ \{2, 1, 3, 4, 5, 6, 7, 8\} \]
\[ \{2, 1, 4, 3, 5, 6, 7, 8\} \]
\[ \{4, 3, 2, 1, 5, 6, 7, 8\} \]
\[ \{4, 3, 2, 1, 6, 5, 7, 8\} \]
\[ \{4, 3, 2, 1, 6, 5, 8, 7\} \]
\[ \{4, 3, 2, 1, 8, 7, 6, 5\} \]
\[ \{8, 7, 6, 5, 4, 3, 2, 1\} \]

\subsection*{(c) Final Array}
The final array is:
\[\{8, 7, 6, 5, 4, 3, 2, 1\}\]

\subsection*{(d) Stack Depth}
The maximum recursion depth is:
\[4\]

\subsection*{(e) Total Recursive Calls}
The total number of recursion calls is:
\[15\]

\newpage
\section*{Question 2}
\subsection*{(a) Differentiation}

\[f(x) = \frac{x^3 - 3x + 2}{\sqrt{x^2 + 1}}\]
Using quotient rule:
\[f'(x) = \frac{u'v - uv'}{v^2}\]
\[u' = 3x^2 - 3, \quad v' = \frac{x}{\sqrt{x^2 + 1}}\]
\[f'(x) = \frac{(3x^2 - 3)\sqrt{x^2 + 1} - (x^3 - 3x + 2)\frac{x}{\sqrt{x^2 + 1}}}{x^2 + 1}\]

\begin{center}
    \fcolorbox{green!50!black}{green!10!white}{\parbox{1\linewidth}{\[ f'(x) = \frac{2x^4 + 3x^2 - 2x - 3}{(x^2 + 1)^{3/2}} \]}}
\end{center}

\subsection*{(b) Integration}

\[\int_1^e x \ln x \, dx\]
Solution:
\[\int x \ln x \, dx = \frac{x^2}{2} \ln x - \frac{x^2}{4} + C\]

\begin{align*}
    \int_1^e x \ln x \, dx &= \left[ \frac{x^2}{2} \ln x - \frac{x^2}{4} \right]_1^e \\
    &= \frac{e^2}{2} - \frac{e^2}{4} + \frac{1}{4}
\end{align*}

\begin{center}
    \fcolorbox{green!50!black}{green!10!white}{\parbox{1\linewidth}{\[ \int_1^e x \ln x \, dx = \frac{e^2 + 1}{4} \]}}
\end{center}

\subsection*{(c) Algebra}

\[2x^2 - 5x + 3 = 0\]
Solution:
\[(2x - 3)(x - 1) = 0\]
\[x = \frac{3}{2}, \quad x = 1\]
Using quadratic formula:
\[x = \frac{5 \pm \sqrt{25 - 24}}{4}\]

\begin{center}
    \fcolorbox{green!50!black}{green!10!white}{\parbox{1\linewidth}{\[ x = \frac{3}{2}, \quad x = 1 \]}}
\end{center}

\newpage
\section*{Question 3}
\subsection*{(a) Determinant and Inverse of Matrix}
Given:
\[A = \begin{pmatrix} 2 & 1 & 0 \\ 1 & 3 & 2 \\ 0 & 2 & 1 \end{pmatrix}\]

Step 1: Determinant
\[\det(A) = 2 \begin{vmatrix} 3 & 2 \\ 2 & 1 \end{vmatrix} - 1 \begin{vmatrix} 1 & 2 \\ 0 & 1 \end{vmatrix} + 0\]
\[= 2(3 \cdot 1 - 2 \cdot 2) - 1(1 \cdot 1 - 2 \cdot 0) = 2(-1) - 1 = -3\]

\begin{center}
    \fcolorbox{blue!60}{blue!5!white}{\parbox{1\linewidth}{\[ \det(A) = -3 \]}}
\end{center}

Step 2: Cofactors (example)
\[C_{11} = -1, \quad C_{12} = -1, \quad C_{13} = 2\]

Step 3: Inverse
\[A^{-1} = \frac{1}{-3} \operatorname{adj}(A)\]

\begin{center}
    \fcolorbox{blue!60}{blue!5!white}{\parbox{1\linewidth}{\[ A^{-1} = -\frac{1}{3} \begin{pmatrix} -1 & -1 & 2 \\ -1 & 2 & -4 \\ 2 & -4 & 5 \end{pmatrix}\]}}
\end{center}

\subsection*{(b) System of Equations (Cramer's Rule)}
\[2x + y = 5\]
\[3x - 2y = 4\]
Solution
\[D = -7, \quad D_x = -14, \quad D_y = -7\]
\[x = \frac{D_x}{D} = 2, \quad y = \frac{D_y}{D} = 1\]

\begin{center}
    \fcolorbox{blue!60}{blue!5!white}{\parbox{1\linewidth}{\[ x = 2, \quad y = 1 \]}}
\end{center}

\subsection*{(c) Vector Operations}

\[\vec{u} = (1, 2, 3), \quad \vec{v} = (4, 5, 6)\]
Solution
\[\vec{u} \cdot \vec{v} = 32\]
\[\vec{u} \times \vec{v} = (-3, 6, -3)\]
\[\theta \approx 12.93^\circ\]

\begin{center}
    \fcolorbox{blue!60}{blue!5!white}{\parbox{1\linewidth}{
    $\text{Final results: Dot } = 32, \text{ Cross } = (-3, 6, -3), \text{ Angle } \approx 12.93^\circ $ }}
\end{center}

\subsection*{(d) Eigenvalues and Eigenvectors}

\[B = \begin{pmatrix} 4 & 1 \\ 2 & 3 \end{pmatrix}\]
Solution
\[\lambda^2 - 7\lambda + 10 = 0 \implies \lambda = 5, 2\]
\[\vec{v}_1 = \begin{pmatrix} 1 \\ 1 \end{pmatrix}, \quad \vec{v}_2 = \begin{pmatrix} 1 \\ -2 \end{pmatrix}\]

\noindent
\fcolorbox{blue!60}{blue!5!white}
    {\parbox{1\linewidth}{
        \text{Eigenvalues: } & 5, 2 \\
        \text{Eigenvectors: } & \begin{pmatrix} 1 \\ 1 \end{pmatrix}, \begin{pmatrix} 1 \\ -2 \end{pmatrix}
}}

\end{document}